\begin{table}[tbp]
\caption{Three completed evaluation designs.}
\label{tab:ch3-designs}
\begingroup\fontsize{9.2}{11.8}\selectfont
\renewcommand{\arraystretch}{1.12}
\begin{tabularx}{\textwidth}{@{}>{\raggedright\arraybackslash}p{23mm}>{\raggedright\arraybackslash}X >{\raggedright\arraybackslash}X@{}}
\toprule
\textbf{Study and question} & \textbf{Implemented evidence} & \textbf{Contribution to Part II}\\
\midrule
Activity: who participates, in what ways, when and where? & Retained GitHub events, a CNCF subset, behavioural aggregation, human/bot distinctions and selected project analyses. & An empirical account of recorded participation. Chapter~4 preserves the definitions, population and 4WH findings.\\[5pt]
OpenRank: what does a relational representation capture? & Formal algorithm, diffusion comparisons, community cases, comparisons among measures and participant assessments. & A relational method examined at several levels. Chapter~5 develops its semantics and evidence; Chapter~6 examines execution.\\[5pt]
ORL: how does evaluation operate when publicly deployed? & A dated leaderboard introduction, project-month series, a common control series, participant survey and interviews. & Evidence about visibility, reported practices and recorded outcomes. Chapter~6 examines deployment and its interpretation.\\
\bottomrule
\end{tabularx}
\endgroup
\tablenote{Reconstructed from \citet{xia2022,han2026rank,zhao2024}. OpenDigger supports observation and execution \citep{wang2026digger}; OpenPerf organizes benchmark roles \citep{han2025openperf}. The OpenRank and ORL accounts share evidence and are not independent replications.}
\end{table}
